\subsection{Real Carbon-Intensity Validation}
\label{sec:evaluation:realcarbon}

To validate that the headline findings of \S\ref{sec:evaluation:headline}
are not artefacts of synthetic carbon traces, we re-ran the canonical
24-hour experiment substituting the synthetic carbon model with the real
ENTSO-E / CAISO 15-minute-resolution carbon-intensity data released by
\citet{wiesner2021lets,lwa}. The workload in \emph{this} experiment
is the schema-correct synthetic stream, holding the workload fixed
while varying only the carbon model; the fully-real version that also
substitutes the real Azure 2021 trace is reported in
\S\ref{sec:evaluation:realboth}.

\paragraph{Carbon profile.} The real LWA traces over our 24-hour
window exhibit mean intensities of 39.9 g/kWh (France), 235.2 g
(Germany), 234.6 g (Great Britain), and 319.0 g (US-CAISO). These are
approximately 15--25\% lower than the calibration values we used in
the synthetic experiments (60, 350, 220, 250 g respectively), reflecting
ENTSO-E's actual 2020 generation mix rather than published annual means.
The diurnal swing is real and non-trivial: GB ranges from 121 to 336
g/kWh on this day, a $\sim 3\times$ ratio; CAISO from 243 to 352, a
$\sim 1.5\times$ ratio. France is essentially flat (38--43 g) due to
its nuclear-dominated mix.

\paragraph{Results.} Table~\ref{tab:real_carbon} reports the same five
schedulers on real carbon data; compare to Table~\ref{tab:headline}.
The qualitative findings of \S\ref{sec:evaluation:headline} hold
without modification.

\begin{table}[h]
\centering\small
\caption{Real LWA carbon, synthetic workload (24h, 173k invocations,
4 regions: DE, FR, GB, US-CAISO),
\textbf{mean $\pm$ std across 5 seeds}. Compare to
Table~\ref{tab:headline}.}
\label{tab:real_carbon}
\begin{tabular}{lrr}
\toprule
Scheduler     & Carbon (g) & vs FIFO \\
\midrule
FIFO          & 33.34 $\pm$ 0.24 & 0.0\% \\
Wait-Awhile   & 92.08 $\pm$ 0.36 & $-$176.2\% $\pm$ 2.5\% \\
Spatial       & 11.63 $\pm$ 0.01 & \textbf{+65.12\% $\pm$ 0.23\%} \\
GreenFaaS-v1  & 14.40 $\pm$ 0.17 & +56.79\% $\pm$ 0.56\% \\
\GreenFaaS    & 11.79 $\pm$ 0.03 & +64.63\% $\pm$ 0.24\% \\
\bottomrule
\end{tabular}
\end{table}

Three observations. \emph{First}, the Wait-Awhile failure mode
reproduces on real data ($-176\% \pm 2.5\%$ over FIFO, vs
$-306\% \pm 4.3\%$ on synthetic). The mechanism is the same
(synchronised-defer pile-up inflating cold-start and idle-energy),
and the magnitude reduction is consistent with the smaller diurnal
swing in the real data.
\emph{Second}, the \GreenFaaS--Spatial gap is $0.49 \pm 0.33$pp,
right on the edge of statistical significance; Spatial slightly edges
\GreenFaaS{}. We attribute this to the absence of a coal-belt region
in the LWA dataset;
under \S\ref{sec:evaluation:topology}, the coal-belt scenario is where
\GreenFaaS{} pulls ahead, and the four-region LWA mix (DE/FR/GB/CAISO)
lacks such a region. \emph{Third}, GreenFaaS-v1's 7pp gap below
\GreenFaaS{} on real data validates the
\S\ref{sec:tradeoff:idle} idle-energy correction in the real-carbon
regime, with the same cold-start-rate increase (0.34\% vs 0.03\%)
seen on synthetic data.

In summary, real carbon data does \emph{not} change any
qualitative finding of \S\ref{sec:evaluation:headline}. The magnitude
of \GreenFaaS's reduction vs FIFO is somewhat smaller on real data
(64\% vs 72\%) because the real European 2020 mix is cleaner overall
than our synthetic calibration assumed. The accompanying topology
sweep on real LWA 2020 traces extended with a calibrated PL trace
(\S\ref{sec:evaluation:realtopology}, Table~\ref{tab:real_topology})
extends this validation across topology variations and tempers the
\S\ref{sec:evaluation:topology} coal-belt claim. The fully-real
validation --- combining the real Azure 2021 per-invocation trace
with time-aligned ElectricityMaps carbon data --- is reported in
\S\ref{sec:evaluation:realboth}, and confirms the qualitative
findings reported here on real workload as well as real carbon.
